\textbf{COLUNI 2025 -- Questão 28}\par\medskip Esqueça o leão! Esqueça o tubarão! Acredita-se que os animais mais letais do mundo não estão na terra e nem no mar, mas “no ar”, são eles, os mosquitos. As doenças transmitidas pelos, aparentemente inofensivos, mosquitos são responsáveis por milhões de mortes a cada ano.

A respeito das doenças transmitidas por mosquitos, julgue as afirmativas abaixo:

I. As doenças são transmitidas, principalmente, por meio da picada do mosquito, mas pode ser também por meio das suas fezes, como no caso, doença de Chagas.\\
II. Entre os principais microrganismos patogênicos transmitidos pelos mosquitos estão os vírus, bactérias, algas e protozoários.\\
III. O mosquito \textit{Aedes aegypti} é o vetor de diversas doenças como febre amarela urbana, febre chikungunya, dengue e amebíase. \\
IV. O uso de telas protetoras em janelas e portas, o uso de repelentes e inseticidas, a eliminação de locais de água parada e a vacinação são algumas formas de controle e prevenção de doenças transmitidas por mosquitos. \\
V. Doenças como a dengue, zika, chikungunya e febre amarela são consideradas doenças negligenciadas uma vez que ocorrem, principalmente, em países pobres, onde há pouco investimento em pesquisas de controle e tratamento.\\

São FALSAS apenas as afirmativas:

\begin{enumerate}
\item II, IV e V.
\item I, III, IV e V.
\item I, II e III.
\item II, III e V.
\end{enumerate}
